\documentclass[11pt]{article}

\usepackage[margin=1in]{geometry}
\usepackage[T1]{fontenc}
\usepackage[utf8]{inputenc}
\usepackage{booktabs}
\usepackage{graphicx}
\usepackage{natbib}
\usepackage{hyperref}

\hypersetup{
	colorlinks=true,
	linkcolor=black,
	citecolor=black,
	urlcolor=blue
}

\title{\bf{Valuing Minnesota Family Forests: \\ Economic and Ecological Contributions}}
\author{Ryan McWay}
\date{\today}

\begin{document}

\maketitle

\begin{abstract}
This report assesses the monetary and non-monetary value of Minnesota family forests, with particular attention to the ecosystem services they provide.
\end{abstract}

\noindent\textbf{Keywords:} Family Forests; Ecosystem Services; Forest Valuation; Minnesota

\section{Introduction}
Describe the importance of Minnesota family forests, the decision context for this assessment, and the research questions addressed by the report.
- Place within the strategic plan of MFRC. 
- Is there a broader (national) context to place this within.

\section{The State of North America's Forests}
Review the relevant literature, regional context, and ecosystem-service concepts that guide the analysis.

- Description of NA, national, and MN forests 
- Narrowing to family forests

\subsection{Minnesota Family Forests}





\section{Methodology to Value Minnesota Family Forests}

\subsection{Study Scope}
Define the study area, forest-owner population, ecosystem services, and time period included in the assessment.


\subsection{Forest Ecosystems Assessed}

- Board overview of ecosystem services 
- Specific ecosystem services evaluated 
- within the context of Gross-ecosystem product (GEP) value and natural capital accounting

\subsubsection{Ecosystem 1:}


\paragraph{Data.}
Describe the datasets, their provenance, preparation, and any important limitations.

\paragraph{Valuation Approach.}
Explain the monetary and non-monetary valuation methods, assumptions, and analytical steps.



\section{What is the Value of Minnesota Family Forests?}

- TBD
- General breakdown of value by groups (which ecosystem services dominate). Comparison of monetary and non-monetary values across ecosystem services. Do market forces drive the effect? 
- Time trends. Is there anything to be stated about changes in forest type, ownership patterns, or ecosystem service provision over time?
- Heterogeneity in forest types and acerage (size)

\section{How Does Forest Value Change under a Warming Climate?}

- Projections of climate change impacts on LULC
- Implications for ecosystem service provision and forest value
- What do we expect for transfer across ownership type (family forest to public or deforestation)?
- What is the implication with agricultural expansion, prarie encroachment, increased global competition driving down timber demand (i.e., Eucalyptus spp.)

\section{Policy Implications}

- Given a fixed MN budget, what policies might we recommend?
- How can policy support the sustainable management and conservation of family forests?
- Are there trade-offs between different policy options in terms of ecosystem service provision and forest value?


\section{Conclusion}
Include limitations here. 
Summarize the principal findings and identify the most important next steps for research, data collection, or decision-making.


\bibliographystyle{plainnat}
\bibliography{references.bib}

\appendix
\section{Supplementary Appendix}
Place detailed methods, assumptions, supporting tables, and additional figures here.

\section{Unassessed Ecosystem Services}
List and describe ecosystem services that were not assessed in this study, explaining why they were excluded and any potential implications for the overall valuation of Minnesota family forests.
Does this lead to underestimation of the total value of Minnesota family forests?

\section{Methodological and Data Limitations}
Here I want to document the limitations related to the methodology, data sources, assumptions, and any uncertainties that may affect the valuation of Minnesota family forests.


\section{Acknowledgments}

\paragraph{Conflicts of Interest.}
The authors declare no conflicts of interest.

\paragraph{Collaborators.}
Support by the Minnesota Family Forest Research Center, and the University of Minnesota Regional Sustainable Development Partnership. 

\paragraph{Funding.}
RSDP funding for Ryan McWay's RA appointment

\end{document}
